Paragraph 2 theme: Capturing blink activity through EEG and EOG signals.
In the neuroscience study of fatigue driving, blink activity can be captured through external methods, such as camera-based eye monitoring, or through electrophysiological recordings, particularly electroencephalography and electrooculography. Although EEG is mainly used to measure brain activity, blink-related components are often present in EEG signals because eyelid movements and eye movements generate strong electrical potentials. These ocular activities are especially prominent in frontal EEG electrodes due to their close proximity to the eyes. As a result, blink information can be extracted from EEG recordings, either by analysing the ocular components embedded in the EEG signal or by combining EEG with EOG-based measurements. This makes EEG/EOG recordings valuable not only for studying neural activity during fatigue but also for identifying blink-related markers associated with drowsiness.